\section{Introduction}
\label{sec:introduction}

\subsection{Motivation}

On 29 October 2024 a DANA — an isolated depression at high levels — struck the
province of Valencia, Spain. It killed 223 people, displaced roughly 15{,}000
residents and caused losses estimated above €50 billion. The Spanish
meteorological agency AEMET issued its highest-level red warning hours before
the worst of the rainfall, but the regional ES-Alert message to citizens'
mobile phones was not sent until late that evening. People died in their
homes, in underground garages, and in vehicles overtaken by water.

The gap that day was not a gap in forecasting. It was a gap between what was
known and what reached the people affected, and between a hazard beginning and
anyone confirming that it had begun. The 2021 eruption of Cumbre Vieja on La
Palma, which forced thousands to evacuate over 85 days, showed the same
pattern from a different hazard.

Spanish cities already carry tens of thousands of cameras: municipal traffic
cameras, road authority cameras, private security systems. During the DANA
those cameras were watching the water rise. Nobody was watching the cameras.

\subsection{What VIGÍA is}

VIGÍA integrates the best publicly available hazard detectors under a single
gated pipeline and adds a shared temporal validation layer that makes them
operationally trustworthy on ordinary municipal cameras.

The system is an integration project and does not claim novel detection
architectures. Every detector it runs was trained by someone else, and each is
credited by name, licence and source in \cref{sec:attribution}. The original
contributions are the validation layer of \cref{sec:validator} and the
evaluation methodology of \cref{sec:methodology}, which measures what that
layer buys and what it costs.

\subsection{The problem the system actually solves}

Hazard detection models are widely available and, individually, good. What
prevents them from being deployed on municipal cameras is not accuracy but
trust: they alarm far too often. Industry reporting puts the proportion of
false alarms in security camera systems at roughly 98\,\%. Commercial wildfire
operators compensate with human staffing — Pano AI charges around
USD 50{,}000 per camera per year and runs a 24/7 intelligence centre, and
ALERTCalifornia routes detections from more than 1{,}060 cameras through
trained staff before dispatch.

We measured this directly rather than citing it. Run exactly as published, the
PyroNear wildfire detector~\citep{pyronear2025} raises an alarm on
\FireImageFAR{} of frames containing no smoke when evaluated against a curated
hard-negative set of \FireNegatives{} images spanning 21 cameras
(\cref{sec:fire}). The detector is not defective; that is the operating point
its authors deliberately chose, assuming a downstream filter. VIGÍA is that
filter.

\subsection{Scope and limitations}

\begin{itemize}
  \item VIGÍA is a research prototype, not a certified emergency system. No
        component has regulatory approval for life-safety use.
  \item It performs no biometric identification of any kind, by architecture
        rather than by configuration (\cref{sec:privacy}).
  \item Each detector carries a declared operating envelope, enforced before
        any model is loaded. Claims outside that envelope are not made.
  \item Several figures come from small samples. Each is reported with the
        interval its sample size supports (\cref{sec:metrics}), and the
        binding cases are listed in \cref{sec:outstanding}.
  \item One hazard — road incidents — is retained as a generality experiment
        for the validator rather than as a disaster capability, and has no
        standalone measurement against held-out labels.
\end{itemize}
